% =============================================================================
\section{Flowcharts}
\label{sec:flow}
% =============================================================================

\tikzset{
  fc/.style={draw, thick, align=center, font=\small, minimum height=0.85cm, fill=white, inner sep=3pt},
  fcproc/.style={fc, rectangle, rounded corners=2pt},
  fcterm/.style={fc, rounded rectangle, fill=gray!15, minimum width=2.4cm},
  fcdec/.style={draw, thick, diamond, aspect=1.6, align=center, font=\footnotesize, inner sep=2pt, fill=stPause!40},
  fa/.style={-{Stealth[length=2.2mm]}, thick},
  fclbl/.style={font=\footnotesize, fill=white, inner sep=1.5pt},
}

\subsection{Main program}

\begin{figure}[H]
  \centering
  \begin{tikzpicture}
    % Left Column: Power-On Initialization Pipeline (centered at x = -4.5)
    \node[fcterm] (start) at (-4.5, 0.0) {Start (\code{Reset_Handler})};
    \node[fcproc, below=0.6cm of start, text width=4.0cm] (clk) 
      {\textbf{Hardware Init}\\[1pt]\footnotesize Core PLL 400\,MHz, Flash, GPIO};
    \node[fcproc, below=0.6cm of clk, text width=4.0cm] (tim) 
      {\textbf{Timer Setup}\\[1pt]\footnotesize SysTick 1.0\,kHz Tick};
    \node[fcproc, below=0.6cm of tim, text width=4.0cm] (fsm) 
      {\textbf{Peripherals Init}\\[1pt]\footnotesize FSM Init, ST7735, UART};
    \node[fcterm, fill=blue!10, below=0.6cm of fsm, text width=3.8cm] (loop_entry) 
      {Enter Super-Loop};

    % Right Area: Non-Blocking Super-Loop Engine
    % Dispatcher decision diamond at (2.5, 0.0)
    \node[fcdec, text width=2.2cm, aspect=1.4] (ms) at (2.5, 0.0) {Elapsed Tick?\\$t_{\text{proc}} < t_{\text{now}}$?};

    % 1ms Tick Catch-Up Branch (Downwards at x = 0.2, width 3.0cm)
    \node[fcproc, text width=3.0cm] (p1) at (0.2, -2.0) 
      {Service 1\,ms Tick\\[1pt]\footnotesize Debounce Integrator};
    \node[fcproc, below=0.55cm of p1, text width=3.0cm] (in) 
      {Sample Inputs\\[1pt]\footnotesize Gesture Decoder};
    \node[fcproc, below=0.55cm of in, text width=3.0cm] (ev) 
      {Dispatch Events\\[1pt]\footnotesize Deterministic EFSM};
    \node[fcdec, below=0.55cm of ev, text width=2.0cm, aspect=1.4] (sec) 
      {1000\,ms\\Elapsed?};
    \node[fcproc, below=0.55cm of sec, text width=3.0cm] (tick) 
      {Dispatch 1\,s Tick\\[1pt]\footnotesize Decrement $t \gets t-1$};

    % Background Idle & Presentation Branch (Downwards at x = 4.8, width 3.0cm)
    \node[fcproc, text width=3.0cm] (uart) at (4.8, -2.0) 
      {Poll USART1\\[1pt]\footnotesize Telemetry / Commands};
    \node[fcproc, below=0.55cm of uart, text width=3.0cm] (out) 
      {Update LEDs\\[1pt]\footnotesize PE3 = $\text{RLED} \lor \text{BLED}$};
    \node[fcproc, below=0.55cm of out, text width=3.0cm] (log) 
      {Stream Telemetry\\[1pt]\footnotesize USART1 (State Log)};
    \node[fcproc, below=0.55cm of log, text width=3.0cm] (draw) 
      {Render Dashboard\\[1pt]\footnotesize ST7735 LCD (If Dirty)};

    % Initialization pipeline connections
    \draw[fa] (start) -- (clk);
    \draw[fa] (clk) -- (tim);
    \draw[fa] (tim) -- (fsm);
    \draw[fa] (fsm) -- (loop_entry);

    % Loop entry connection routing cleanly around open corridor
    \draw[thick] (loop_entry.east) -- ++(1.0,0) |- (2.5, 1.8);

    % Super-loop branch connections
    \draw[fa] (ms.west) -| node[fclbl, above, pos=0.35]{yes (catch-up)} (p1);
    \draw[fa] (p1) -- (in);
    \draw[fa] (in) -- (ev);
    \draw[fa] (ev) -- (sec);
    \draw[fa] (sec) -- node[fclbl, right=2pt, pos=0.45]{yes} (tick);
    
    \draw[fa] (ms.east) -| node[fclbl, above, pos=0.35]{no (idle)} (uart);
    \draw[fa] (uart) -- (out);
    \draw[fa] (out) -- (log);
    \draw[fa] (log) -- (draw);

    % Loop-back path merge (meeting at corridor x = 1.8 below tick.south)
    \coordinate (merge) at ($(tick.south) + (1.6cm, -0.8cm)$);
    \draw[thick] (tick.south) |- (merge);
    \draw[thick] (sec.east) -- node[fclbl, above, pos=0.3]{no} ++(0.6,0) |- (merge);
    \draw[thick] (draw.south) |- (merge);
    \draw[thick] (merge) -- ++(5.35,0) |- (2.5, 1.8);
    \draw[fa] (2.5, 1.8) -- (ms.north);
  \end{tikzpicture}
  \caption{Main program flowchart: Non-blocking super-loop architecture}
  \label{fig:flow-main}
\end{figure}

\subsection{Coin processing}

\begin{figure}[H]
  \centering
  \begin{tikzpicture}[node distance=0.7cm]
    % Discrete columns: x = 2.8 (Idle branch), x = 8.0 (Ready branch), x = 12.2 (Reject branch)
    \node[fcterm] (s) {COIN($v$) Event Received};
    \node[fcdec, below=0.8cm of s, text width=2.4cm] (idle) {State Idle?\\(\st{STANDBY} / \st{COLLECT})};
    \node[fcdec, right=1.6cm of idle, text width=2.2cm] (rdy) {State =\\ \st{READY}?};
    \node[fcterm, fill=gray!15, right=1.2cm of rdy] (rej) {Ignore Coin\\[2pt]\footnotesize Active / Error};

    \node[fcproc, below=0.9cm of idle, text width=2.8cm] (add) {$b \gets b + v$\\[2pt]\footnotesize (Guarded sum)};
    \node[fcproc, below=0.9cm of rdy, text width=2.8cm] (add2) {$b \gets b + v$\\[2pt]\footnotesize (Surplus deposit)};

    \node[fcdec, below=0.9cm of add, text width=2.0cm] (th) {$b \ge 50\,\text{\textcent}$?};
    \node[fcproc, below left=0.9cm and 0.2cm of th, text width=2.6cm] (col) {Enter \st{COLLECTING}\\[2pt]\footnotesize RLED ON, BLED OFF};
    \node[fcproc, below right=0.9cm and 0.2cm of th, text width=2.6cm] (ready) {Enter \st{READY}\\[2pt]\footnotesize RLED OFF, BLED ON};
    \node[fcterm, below=3.4cm of th, text width=3.4cm] (e) {Return / End Event};

    \draw[fa] (s) -- (idle);
    \draw[fa] (idle) -- node[fclbl, right=2pt, pos=0.45]{yes} (add);
    \draw[fa] (idle) -- node[fclbl, above=2pt]{no} (rdy);
    \draw[fa] (rdy) -- node[fclbl, above=2pt]{no} (rej);
    \draw[fa] (rdy) -- node[fclbl, right=2pt, pos=0.45]{yes} (add2);
    \draw[fa] (add) -- (th);
    \draw[fa] (th) -| node[fclbl, above, pos=0.3]{no} (col);
    \draw[fa] (th) -| node[fclbl, above, pos=0.3]{yes} (ready);
    \draw[fa] (col.south) |- (e.west);
    \draw[fa] (ready.south) -- (ready.south |- e.north);
    \draw[thick] (rej.south) |- (add2.south);
    \draw[fa] (add2.south) |- (e.east);
  \end{tikzpicture}
  \caption{Coin processing logic in the deterministic FSM core}
  \label{fig:flow-coin}
\end{figure}

\noindent On the board, coins come from K1 gestures and are first \emph{staged} by the front panel
(Figure~\ref{fig:flow-staging}); a real coin acceptor would instead feed the hardware pulse accumulator driver, which
emits COIN($v$) events directly.

\begin{figure}[H]
  \centering
  \begin{tikzpicture}[node distance=0.7cm]
    \node[fcterm] (s) {K1 Gesture in \st{STANDBY}};
    \node[fcdec, below=0.8cm of s, text width=2.0cm] (g) {Gesture\\Type?};
    \node[fcproc, left=1.8cm of g, text width=3.8cm] (add) {Pending $\gets$ Pending $+ v$\\[2pt]\footnotesize $10\,\text{\textcent}, 20\,\text{\textcent}, 50\,\text{\textcent}$ (Click/Double/Hold)\\[1pt]\footnotesize Guarded Max $990\,\text{\textcent}$};
    \node[fcdec, below=1.2cm of g, text width=2.0cm] (th) {Pending\\$\ge 50\,\text{\textcent}$?};
    \node[fcproc, below left=1.2cm and 0.2cm of th, text width=3.6cm] (ok) {Deposit $50\,\text{\textcent}$ to FSM Core\\[2pt]\footnotesize Advance to \st{READY}\\[1pt]\footnotesize Display: ``Change Pending$-$50''};
    \node[fcproc, below right=1.2cm and 0.2cm of th, text width=3.6cm] (no) {Remain in \st{STANDBY}\\[2pt]\footnotesize Display: ``Not Enough, Refund''};
    \node[fcproc, below=4.2cm of th, text width=3.6cm] (clr) {Clear Staged Balance\\[2pt]\footnotesize Pending $\gets 0$};
    \node[fcterm, below=0.7cm of clr] (e) {Wait for Next Event};

    \draw[fa] (s) -- (g);
    \draw[fa] (g) -- node[fclbl, above=2pt]{Deposit} (add);
    \draw[fa] (add.north) |- ($(s.south)!0.5!(g.north)$);
    \draw[fa] (g) -- node[fclbl, right=2pt, pos=0.45]{HOLD2 (1.2\,s) = Confirm} (th);
    \draw[fa] (th) -| node[fclbl, above, pos=0.3]{yes} (ok);
    \draw[fa] (th) -| node[fclbl, above, pos=0.3]{no} (no);
    \draw[fa] (ok.south) |- (clr.west);
    \draw[fa] (no.south) |- (clr.east);
    \draw[fa] (clr) -- (e);
  \end{tikzpicture}
  \caption{Coin staging logic on the board single-key front panel}
  \label{fig:flow-staging}
\end{figure}

\subsection{RUN / PAUSE / STOP handling}

\begin{figure}[H]
  \centering
  \resizebox{\linewidth}{!}{%
  \begin{tikzpicture}
    % Top entry and main button discriminator (relative placement, zero inversion)
    \node[fcterm] (s) at (0, 0.0) {User Button Event Received};
    \node[fcdec, below=0.8cm of s, text width=2.0cm, aspect=1.4] (k) {Which\\Button?};

    % Column 2: PAUSE logic (directly below k)
    \node[fcdec, below=1.0cm of k, text width=2.0cm, aspect=1.4] (p1) {State =\\\st{RUNNING}?};
    \node[fcproc, below=0.75cm of p1, text width=3.0cm] (p1_act) 
      {\textbf{Pause Cycle}\\[2pt]\footnotesize Motor OFF, BLED Solid ON\\Door Remains LOCKED\\$\rightarrow$ \st{PAUSED} ($t$ Continues)};
    \node[fcterm, below=0.75cm of p1_act, text width=2.2cm] (p_ign) {Ignore PAUSE};

    % Column 1: RUN logic (left of p1)
    \node[fcdec, left=2.4cm of p1, text width=2.0cm, aspect=1.4] (r1) {State =\\\st{READY}?};
    \node[fcproc, below=0.65cm of r1, text width=3.0cm] (r1_act) 
      {\textbf{Start Cycle}\\[2pt]\footnotesize $b \gets 0, \; t \gets 1800\,\text{s}$,\\Door LOCKED, Motor ON\\$\rightarrow$ \st{RUNNING}};
    \node[fcdec, below=0.65cm of r1_act, text width=2.0cm, aspect=1.4] (r2) {State =\\\st{PAUSED}?};
    \node[fcproc, below=0.65cm of r2, text width=3.0cm] (r2_act) 
      {\textbf{Resume Cycle}\\[2pt]\footnotesize Motor ON, BLED 1\,Hz\\$t$ Preserved\\$\rightarrow$ \st{RUNNING}};
    \node[fcterm, below=0.65cm of r2_act, text width=2.2cm] (r_ign) {Ignore RUN};

    % Column 3: STOP logic (right of p1)
    \node[fcdec, right=2.4cm of p1, text width=2.0cm, aspect=1.4] (s1) {\st{STANDBY} or\\\st{ERROR}?};
    \node[fcterm, below=0.65cm of s1, text width=2.2cm] (s_ign) {Ignore STOP};
    \node[fcdec, below=0.65cm of s_ign, text width=2.0cm, aspect=1.4] (s2) {Double-Press\\$n = 1$?};
    \node[fcproc, below=0.65cm of s2, text width=3.2cm] (s2_arm) 
      {\textbf{Arm Window}\\[2pt]\footnotesize $n \gets 1, \; w \gets 1.5\,\text{s}$};
    \node[fcproc, below=0.65cm of s2_arm, text width=3.2cm] (s2_abort) 
      {\textbf{Force Stop / Reset}\\[2pt]\footnotesize Refund deposit $b$\\Actuators OFF, Unlock Door\\$\rightarrow$ \st{STANDBY}};

    % Main distribution routing
    \draw[fa] (s) -- (k);
    \draw[fa] (k) -| node[fclbl, above, pos=0.35]{RUN} (r1);
    \draw[fa] (k) -- node[fclbl, right=2pt, pos=0.45]{PAUSE} (p1);
    \draw[fa] (k) -| node[fclbl, above, pos=0.35]{STOP} (s1);

    % RUN branch routing
    \draw[fa] (r1) -- node[fclbl, right=2pt, pos=0.45]{yes} (r1_act);
    \draw[fa] (r1.west) -- ++(-0.8,0) node[fclbl, midway, above]{no} |- (r2.west);
    \draw[fa] (r2) -- node[fclbl, right=2pt, pos=0.45]{yes} (r2_act);
    \draw[fa] (r2.east) -- ++(0.8,0) node[fclbl, midway, above]{no} |- (r_ign.east);

    % PAUSE branch routing
    \draw[fa] (p1) -- node[fclbl, right=2pt, pos=0.45]{yes} (p1_act);
    \draw[fa] (p1.east) -- ++(0.9,0) node[fclbl, midway, above]{no} |- (p_ign.east);

    % STOP branch routing
    \draw[fa] (s1) -- node[fclbl, right=2pt, pos=0.45]{yes} (s_ign);
    \draw[fa] (s1.west) -- ++(-0.8,0) node[fclbl, midway, above]{no} |- (s2.west);
    \draw[fa] (s2) -- node[fclbl, right=2pt, pos=0.45]{no} (s2_arm);
    \draw[fa] (s2.east) -- ++(1.4,0) node[fclbl, midway, above]{yes ($w > 0$)} |- (s2_abort.east);
  \end{tikzpicture}%
  }
  \caption{Button gesture handling logic across discrete operational modes}
  \label{fig:flow-buttons}
\end{figure}

\subsection{Multi-Gesture Decoding Engine}

\begin{figure}[H]
  \centering
  \begin{tikzpicture}
    % Top entry (relative placement, zero inversion)
    \node[fcterm] (s) at (0, 0.0) {Debounced Key State Sampled (1\,ms)};
    \node[fcdec, below=0.85cm of s, text width=2.0cm, aspect=1.4] (down) {Key Held\\DOWN?};

    % Left Column: Hold Duration Discrimination
    \node[fcproc, below left=0.9cm and 1.2cm of down, text width=3.2cm] (d1) 
      {State $\gets$ \code{DOWN1}\\[2pt]\footnotesize $t_{\text{down}} \gets t_{\text{down}} + 1\,\text{ms}$};
    \node[fcdec, below=0.65cm of d1, text width=2.0cm, aspect=1.4] (h1) {$t_{\text{down}} \ge 600\,\text{ms}$?};
    \node[fcproc, below=0.65cm of h1, text width=3.2cm] (hold1) 
      {Emit \code{HOLD1}\\[2pt]\footnotesize STOP \#1 / 50\cent{} Coin};
    \node[fcdec, below=0.65cm of hold1, text width=2.0cm, aspect=1.4] (h2) {$t_{\text{down}} \ge 1200\,\text{ms}$?};
    \node[fcproc, below=0.65cm of h2, text width=3.2cm] (hold2) 
      {Emit \code{HOLD2}\\[2pt]\footnotesize STOP \#2 / Confirm Fare};

    % Right Column: Multi-Click Timing Discrimination
    \node[fcproc, below right=0.9cm and 1.2cm of down, text width=3.2cm] (wait) 
      {Key Released\\State $\gets$ \code{WAIT}\\[2pt]\footnotesize $t_{\text{gap}} \gets t_{\text{gap}} + 1\,\text{ms}$};
    \node[fcdec, below=0.65cm of wait, text width=2.0cm, aspect=1.4] (gap) {$t_{\text{gap}} > 250\,\text{ms}$?};
    \node[fcproc, below=0.65cm of gap, text width=3.0cm] (clk) 
      {Emit \code{CLICK}\\[2pt]\footnotesize State $\gets$ \code{IDLE}};
    \node[fcdec, below=0.65cm of clk, text width=2.0cm, aspect=1.4] (d2) {Second Press\\within $250\,\text{ms}$?};
    \node[fcproc, below=0.65cm of d2, text width=3.0cm] (dbl) 
      {Emit \code{DOUBLE}\\[2pt]\footnotesize State $\gets$ \code{IDLE}};

    % Routing
    \draw[fa] (s) -- (down);
    \draw[fa] (down) -| node[fclbl, above, pos=0.35]{yes} (d1);
    \draw[fa] (down) -| node[fclbl, above, pos=0.35]{released} (wait);

    \draw[fa] (d1) -- (h1);
    \draw[fa] (h1) -- node[fclbl, right=2pt, pos=0.45]{yes} (hold1);
    \draw[fa] (hold1) -- (h2);
    \draw[fa] (h2) -- node[fclbl, right=2pt, pos=0.45]{yes} (hold2);

    \draw[fa] (wait) -- (gap);
    \draw[fa] (gap) -- node[fclbl, right=2pt, pos=0.45]{yes (timeout)} (clk);
    \draw[fa] (gap.east) -- ++(0.9,0) node[fclbl, midway, above]{no} |- (d2.east);
    \draw[fa] (d2) -- node[fclbl, right=2pt, pos=0.45]{yes} (dbl);
  \end{tikzpicture}
  \caption{Multi-gesture button decoder sub-machine flowchart}
  \label{fig:flow-gesture}
\end{figure}

\subsection{Emergency Safety Fault Handling and Context Recovery}

Figure~\ref{fig:flow-fault} illustrates the fail-safe emergency shutdown procedure and context-restoring
recovery path governed by Safety Invariant $I_3$ and Assumption A8.

\begin{figure}[H]
  \centering
  \begin{tikzpicture}[node distance=0.65cm]
    \node[fcterm, fill=red!15] (fault) {Safety Fault Asserted};
    \node[fcproc, below=of fault, text width=9.2cm] (save) 
      {\textbf{Snapshot Active Operational State}: $X_{\text{saved}} \gets S$\\[2pt]\footnotesize Record Active Fault Flags: $f \gets \text{flags}$};
    \node[fcproc, below=of save, text width=9.2cm] (safe) 
      {\textbf{Immediate Fail-Safe Isolation} ($T_{\text{react}} \le 30\,\text{ms}$):\\[2pt]\footnotesize De-energize Actuators (Motor, Pump, Valve), Release Door Deadbolt,\\[1pt]\footnotesize Activate 2.0\,Hz RLED Optical Alarm};
    \node[fcproc, below=of safe, text width=9.2cm] (enter) 
      {\textbf{Enter \st{ERROR} Safe State}\\[2pt]\footnotesize Lock Out All User Control Command Inputs};
    \node[fcdec, below=of enter, text width=2.6cm, aspect=1.4] (wait) {Fault Cleared\\or Reset?};
    \node[fcdec, below=1.0cm of wait, text width=3.0cm, aspect=1.3] (was_cycle) {$X_{\text{saved}} \in$\\$\{\text{\st{RUNNING}}, \text{\st{PAUSED}}\}$\\$\land\; t > 0$?};
    \node[fcproc, below left=1.0cm and 0.4cm of was_cycle, text width=4.2cm] (res_pause) 
      {\textbf{Transition to \st{PAUSED}}\\[2pt]\footnotesize Actuators Stay Suspended\\BLED Solid ON\\Door LOCKED\\Awaiting RUN to Resume Cycle};
    \node[fcproc, below right=1.0cm and 0.4cm of was_cycle, text width=4.2cm] (res_idle) 
      {\textbf{Restore Idle State by Balance}:\\[2pt]\footnotesize $b \ge 50 \implies \text{\st{READY}}$\\$0 < b < 50 \implies \text{\st{COLLECTING}}$\\$b = 0 \implies \text{\st{STANDBY}}$};

    \draw[fa] (fault) -- (save);
    \draw[fa] (save) -- (safe);
    \draw[fa] (safe) -- (enter);
    \draw[fa] (enter) -- (wait);
    \draw[fa] (wait) -- node[fclbl, right=2pt, pos=0.45]{yes} (was_cycle);
    \draw[fa] (wait.east) -- ++(4.6,0) |- node[pos=0.25, right, fclbl]{no ($1\,\text{s}$ Tick active)} (enter.east);
    \draw[fa] (was_cycle) -| node[fclbl, above, pos=0.3]{yes} (res_pause);
    \draw[fa] (was_cycle) -| node[fclbl, above, pos=0.3]{no} (res_idle);
  \end{tikzpicture}
  \caption{Emergency fault trip and context-preserving recovery flowchart}
  \label{fig:flow-fault}
\end{figure}
